\documentclass[a4paper,12pt]{article}
\usepackage[portuges]{babel}
\usepackage[latin1]{inputenc}
\usepackage{graphicx}
\usepackage{amsmath,amsfonts,mathrsfs,amssymb}
\usepackage{multicol}
\usepackage{fancyhdr}
\usepackage{makeidx}
\numberwithin{equation}{section}
\usepackage{geometry}
\geometry{verbose,tmargin=2cm,bmargin=1cm,lmargin=2cm,rmargin=2cm}
\usepackage[pdftex,plainpages=false,pdfpagelabels,pagebackref,colorlinks=true,
citecolor=black,linkcolor=black,urlcolor=black,filecolor=black,bookmarksopen=true]{
hyperref}
\usepackage{sectsty}
\sectionfont{\large}
\subsectionfont{\normalsize}

\setlength{\textheight}{240mm}
\setlength{\headheight}{28pt}
\renewcommand{\headrulewidth}{0.5pt}
\renewcommand{\footrulewidth}{0.5pt}
\renewcommand{\baselinestretch}{1.2}

%Atalhos para nota��o de conjuntos
\newcommand{\cqd}{{{\hfill $\rule{2.5mm}{2.5mm}$}}\vspace{1cm}\\}
\newcommand{\sen}{\mathrm{sen}}
\newcommand{\ds}{\displaystyle}
\newcommand{\R}{\mathbb{R}}
\newcommand{\C}{\mathbb{C}}
\newcommand{\N}{\mathbb{N}}
\newcommand{\Z}{\mathbb{Z}}
\newcommand{\Q}{\mathbb{Q}}
\newcommand{\I}{\mathbb{I}}

%Comandos Atalhos para opera��es matem�ticas
\newcommand{\deli}[2]{\ensuremath\frac{\partial#1}{\partial#2}}
\newcommand{\deri}[2]{\ensuremath\frac{d#1}{d#2}}
\newcommand{\nt}[4]{\ensuremath{ \int_{#1}^{#2} #3 \, \mbox{d}#4}}

%Atalhos para nota��es matem�ticas, fun��es, conjuntos, etc.
\newcommand{\raiotres}{\sqrt{x^2 + y^2 + z^2}}
\newcommand{\coordn}[1]{(#1_1,\dotsc,#1_n)}
\newcommand{\setn}[1]{(#1_1,\dotsc,#1_n)}
\newcommand{\prodintn}[2]{#1_1#2_1 + \dotsc + #1_n#2_n}
\newcommand{\traco}{{\text{Tr}}}
\newcommand{\posto}{{\text{posto}}}
\newcommand{\Ima}{{\text{Im}}}

%Atalhos para ambientes matem�ticos
\newcommand{\exercicio}[1]{\vspace{0.2cm}\noindent\textbf{Exerc�cio #1. }\ }
\newcommand{\propriedade}[1]{\vspace{0.2cm}\noindent\textbf{Propriedade #1:}\ }
\newcommand{\definicao}{\vspace{0.2cm}\noindent\textbf{Defini��o:}\ }
\newcommand{\solucao}{\vspace{0.2cm}\noindent\textbf{Solu��o:}\ }
\newcommand{\exemplo}[1]{\vspace{0.2cm}\noindent\textbf{Exemplo #1 :}\ }
\newcommand{\demonstracao}{\vspace{0.2cm}\noindent\textbf{Demonstra��o:}\ }

%Atalhos para nota��es vetoriais
\newcommand{\diff}{\text{d}}
\newcommand{\te}{\text{t}}
\newcommand{\di}{\mbox{d}}
\newcommand{\x}{{\bf x}}
\newcommand{\versor}[1]{{\bf #1}}
\newcommand{\y}{{\bf y}}
\newcommand{\e}{{\bf e}}
\newcommand{\ve}{{\bf v}}

\pagestyle{fancy}
\fancyhead[LO,L]{\textit \normalsize \nouppercase \leftmark}
\fancyhead[RO,R]{\textit \normalsize \nouppercase \rightmark}
\fancyfoot[LO,L]{\itshape{Osvaldo Pereira}}
\fancyfoot[RO,R]{Equa��o de Dirac}

\begin{document}
\tableofcontents
\section{Motiva��o Hist�rica}

\subsection{Introdu��o}

Depois que a teoria da Relatividade Especial \index{Relatividade Especial} foi
estabelecida definitivamente no meio cient�fico, contempor�neamente come�ou-se a formular
a teoria qu�ntica. O advendo da relatividade incentivou diversos cientistas a procurarem
formular a mec�nica qu�ntica \index{mec�nica qu�ntica} de uma forma (Lorentz)
\index{Lorentz} covariante \index{covariante}, da mesma maneira que aconteceu com o
eletromagnetismo \index{eletromagnetismo} de Maxwell.
\\
Na �poca, na d�cada de 1920, as �nicas part�culas conhecidas eram el�tron \index{el�tron}
, pr�ton \index{pr�ton} e f�ton \index{f�ton} . Ainda n�o se falava em b�sons
\index{b�sons} , f�rmions \index{f�rmions} , e logicamente em estat�stica de part�culas,
t�o pouco era conhecido ainda o conceito de spin \index{spin} de uma part�cula. Mas j� se
tentava descrever a din�mica \index{din�mica} dessas part�culas, sob a a��o de um
potencial escalar. Emprestava-se os conceitos de colchetes de Poisson \index{colchetes de
Poisson}, no qual descrevia-se classicamente a din�mica de part�culas, e foi Schrodinger,
com sua equa��o diferencial \index{equa��o diferencial}, quem primeiro conseguiu
fundamentar uma descri��o puramente qu�ntica.
%
%
\[
i\hbar\frac{\partial}{\partial t} \Psi(\mathbf{r},\,t) =
-\frac{\hbar^2}{2m}\nabla^2\Psi(\mathbf{r},\,t) + V(\mathbf{r})\Psi(\mathbf{r},\,t)
\]
%
%
A equa��o acima nada mais � do que uma equa��o de auto valores \index{equa��o de auto
valores}, que descreve a evolu��o temporal \index{evolu��o temporal} de uma fun��o (auto
fun��o), $\psi(\versor{x},t)$, de forma que:
%
%
\[H\psi(\versor{x},t) = i\hbar\deli{\psi}{t}(\versor{x},t)\]
%
%
Onde H � a fun��o Hamiltoniana \index{Hamiltoniana} (energia total do sistema), dada
por:
%
%
\[H = \frac{p^2}{2m} + V(\versor{x})\]
%
%
Onde a primeira parcela do lado direito nada mais � do que a energia cin�tica
\index{energia cin�tica} da part�cula de massa m, e a segunda parcela seria o potencial
escalar \index{potencial escalar}, sob o qual a part�cula sofre intera��o
\index{intera��o}.
\\
Depois que a teoria qu�ntica foi estabelecida no meio cient�fico, credita-se o
sucesso dessa equa��o � passagem da mec�nica cl�ssica para a qu�ntica, tratando o momento
linear, e a posi��o n�o apenas mais como fun��es que descreviam univocamente a trajet�ria
cl�ssica de uma part�cula, mas agora como operadores. E a energia do sistema
\index{energia do sistema} $(\epsilon)$ seria a auto solu��o \index{auto solu��o} de uma
equa��o de auto valores:
%
%
\[H\psi = \epsilon\psi\]
%
%
A transi��o da cl�ssica para o qu�ntico, foi feita da seguinte maneira:
%
%
\[p\to i\hbar\nabla \qquad\qquad E \to i\hbar\deli{}{t}\]
%
%
O que motivou Schrodinger a descobrir sua equa��o foi a an�lise qualitativa e quantitativa
de espectros at�micos, primeiramente descritos com certa fidelidade, a partir do Modelo de
Bohr \index{Modelo de Bohr}. A equa��o de Schrodinger descreve a din�mica de uma
part�cula, com spin 0, em um potencial escalar.
\\
Heisenberg \index{Heisenberg} que tamb�m trabalhava na descri��o din�mica de part�culas,
conseguiu fazer a mesma an�lise que Schrodinger, por�m em um outro enfoque. Na descri��o
de Schrodinger, os operadores momento e posi��o eram tratados como observadores, e a
evolu��o temporal era toda realizada nas auto fun��es, enquanto que na descri��o de
Heisenberg os observadores eram evoluidos unitariamente no tempo, e as auto fun��es eram
solu��es estacion�rias. Mais tarde comprovou-se que as duas descri��es eram
complementares, atrav�s de uma transforma��o de similaridade.
\\
Schrodinger descobriu sua equa��o em 1926, e o entendimento desta descri��o qu�ntica da
din�mica de part�culas, em formato n�o relativ�stico, durou cerca de dois anos. Enquanto
que o entendimento da sua formula��o covariante (relativ�stica) demorou cerca de seis
anos. Surgia ent�o a teoria qu�ntica relativ�stica.
\\
A equa��o de Dirac surgiu em 1928, mas paralelamente � Dirac, tanto Schrodinger quanto
Pauli, Klein e Gordon, trabalhavam em problemas similares. H� evid�ncias inclusive de que
o pr�prio Schrodinger haveria descoberto uma formula��o relativ�stica, e teria encontrado
os valores negativos de energia, que Klein e Gordon haveriam encontrado, e que depois
seria magistralmente explicada por Dirac, com sua teoria do par el�tron buracos. Por se
tratar na �poca de algo inusitado, Schrodinger teria desistido de perseguir essa
formula��o, pois os valores negativos de energia n�o teriam significado f�sico.
\\
A equa��o de Pauli, descoberta em 1927, seria a descri��o de part�culas com spin $1/2$.
Foi Pauli quem cunhou o termo spin. Pelo fato de sua equa��o possuir a intera��o das
part�culas com um campo vetorial, sua solu��o n�o seria mais uma fun��o escalar, e sim um
vetor, Pauli nomeou estes vetores de spinores.
%
%
\[
\left[ \frac{1}{2m}(\versor{\sigma}\cdot(\versor{p} - q \versor{A}))^2 + q \phi \right]
\psi = i \hbar \frac{\partial}{\partial t}\psi
\]
%
%
Onde m � a massa da part�cula, $\phi$ o potencial escalar, $\versor{A}$ seria o potencial
vetor, e $\versor{\sigma}$ seria um 3-vetor, cujas compontentes seriam as matrizes que
Pauli descobrira.


\subsection{Equa��o de Klein Gordon}

A equa��o de Schrodinger � uma equa��o de auto valores, $H\psi = E\psi$, onde $H$ � a
fun��o Hamiltoniana do sistema, sendo a energia total:
%
%
\begin{eqnarray}
\nonumber H\psi &=& E\psi\\
\nonumber \left(\frac{p^2}{2m} + V\right)\psi &=& E\psi\\
\nonumber  \left(- \frac{h^2}{2m}\nabla^2 + V\right)\psi &=& i\hbar\deli{\psi}{t}
\end{eqnarray}
%
%
Tomando o 4-momento, $p^\mu = i\hbar\partial_\mu$, como sendo da seguinte forma:
%
%
\begin{eqnarray}
\nonumber p^\mu &=& i\hbar (\partial_{ct},- \partial_x,- \partial_y,- \partial_z)\\
\nonumber       &=& i\hbar (\partial_{ct},- \nabla)\\
\nonumber       &=& i\hbar (p_0,\versor{p})
\end{eqnarray}
%
%
� poss�vel obter a equa��o de Klein-Gordon \index{equa��o de Klein-Gordon} para part�culas
livres, na forma de Schrodinger relativ�stica. De maneira que seja tamb�m uma equa��o de
auto valores, $p^\mu p_\mu\psi = (mc)^2\psi$.
%
%
\[p^\mu p_\mu\psi = \left(\frac{E^2}{c^2} - p^2\right)\psi = - (mc)^2\]
%
%
De forma que seus poss�veis auto-valores sejam dados por:
%
%
\[E_+ = \sqrt{(pc)^2 + (mc^2)^2} \qquad , \qquad E_- = - \sqrt{(pc)^2 + (mc^2)^2}\]
%
%
Esses dois poss�veis valores de auto energias na equua��o de Klein Gordon, � a principal
diferen�a, em compara��o com a equa��o de Schrodinger. Substituindo, $E \to
i\hbar\partial_{ct}$, e $p \to i\hbar\nabla$, tem-se que:
%
%
\[- \hbar^2\left(\deli{^2}{(ct)^2} - \nabla^2\right)\psi = - (mc)^2\psi\]
%
%
Como $p^\mu p_\mu$ � invariante de Lorentz \index{invariante de Lorentz}, isso torna a
equa��o de Klein-Gordon covariante \index{Klein-Gordon covariante}. Definindo o seguinte
operador:
%
%
\[\partial^2 = \deli{^2}{(ct)^2} - \nabla^2\]
%
%
Encontra-se a seguinte rela��o (Equa��o de Klein Gordon):
%
%
\[\left(\partial^2 + \frac{m^2c^2}{\hbar^2}\right)\psi = 0\]
%
%
A rela��o acima � uma Equa��o de Poisson \index{Equa��o de Poisson} do tipo, $(\triangle
- \lambda^2)\psi = - f$ . A rela��o $\lambda_c = \hbar/mc$ � o comprimento de onda
Compton \index{comprimento de onda Compton} da part�cula em quest�o. A solu��o da
part�cula livre para a Equa��o de Klein Gordon � a mesma que para a Equa��o
de Schrodinger.
%
%
\[\psi = e^{- ip_\mu x^\mu/\hbar} = e^{ i(\versor{p}\cdot\versor{x} - Et)/\hbar}\]
%
%
� uma consequ�ncia direta da Equa��o de Klein Gordon, que a fun��o de onda satisfaz a
equa��o de continuidade. Tomando a equa��o e seu conjugado, e multiplicando pela esquerda,
pela fun��o de onda e pelo seu complexo conjugado, respectivamente, tem-se que:
%
%
\begin{eqnarray}
\nonumber \psi^*( p_\mu p^\mu - m^2c^2)\psi &=& 0\\
\nonumber \psi( p_\mu p^\mu  - m^2c^2)\psi^* &=& 0
\end{eqnarray}
%
%
Subtraindo uma equa��o da outra, tem-se que:
%
%
\[\psi^*( p_\mu p^\mu - m^2c^2)\psi  - \psi( p_\mu p^\mu - m^2c^2)\psi^*\]
%
%
Tendo-se em mente que $\psi^*\psi = \psi\psi^* = |\psi|^2$, a rela��o acima simplifica
para:
%
%
\[ \psi^* p_\mu p^\mu\psi - \psi p_\mu p^\mu\psi^*  = 0 \]
%
%
Somando e subtraindo $p_\mu\psi^* p^\mu\psi$, e tendo em mente que, $p_\mu\psi^*
p^\mu\psi = p_\mu\psi p^\mu\psi^*$
%
%
\[ \psi^* p_\mu p^\mu\psi + p_\mu\psi^* p^\mu\psi - p_\mu\psi p^\mu\psi^* - \psi p_\mu
p^\mu\psi^*  = 0 \]
%
%
Simplificando, tem-se que:
%
%
\[p_\mu(\psi^* p^\mu\psi) - p_\mu(\psi p^\mu\psi^*) = p_\mu(\psi^* p^\mu\psi - \psi
p^\mu\psi^*) = 0 \]
%
%
Definindo, a 4-corrente de densidade $j^\mu = \psi^* p^\mu\psi - \psi
p^\mu\psi^*$, tem-se a equa��o de continuidade \index{equa��o de continuidade}, $p_\mu
j^\mu = 0$ (em forma tensorial). Em forma vetorial, lembrando que, $p_\mu =
g_{\mu\nu}p^\nu = i\hbar(\partial_{ct},\nabla)$, tem-se que:
%
%
\[\deli{\rho}{t} + \nabla\cdot\versor{j} = 0\]
%
%
Onde definiu-se que:
%
%
\[\rho = \psi^*\deli{}{t}\psi - \psi\deli{}{t}\psi^* \quad , \quad
\versor{j} = \psi\nabla\psi^* - \psi^*\nabla\psi
\]

\subsection{Formula��o Covariante}

Dirac teve como motiva��o para deduzir sua equa��o, deixar a equa��o na forma de
Schrodinger, invariante sobre transforma��es de Lorentz. Podendo ent�o, ser aplicada �
part�culas com velocidades consider�veis, pr�ximas da velocidade da luz.
\\
O grade desafio de Dirac foi mesclar conceitos da mec�nica qu�ntica com os da
relatividade especial. De um lado, da qu�ntica, tinha-se os conceitos de fun��es de
estado, o princ�pio fundamental da superposi��o de estados, e o de observadores f�sicos.
Do outro lado, na relatividade especial, os conceitos de "fun��es de estado", e o conceito
de simultaneidade.
\\
Os conceitos de fun��o de estado na qu�ntica mesclariam de forma precisa com os da
relatividade especial, a fun��o de onda, solu��o da equa��o de Schrodinger,
$\psi(\versor{x},t)$ calharia perfeitamente com a descri��o de um 4-vetor, na
relatividade especial, sendo a fun��o de onda, $\psi = \psi(x_0,x_1,x_2,x_3)$, onde
tinha-se que, $x_0 = ct$, e $\versor{x} = (x_1,x_2,x_3)$ seriam as componentes espaciais
da fun��o de onda.
\\
O grande problema nesse tratamento, foi o conceito geral de um observador. Como na
mec�nica qu�ntica � poss�vel que existam mais de um observador realizando mesmas medidas,
de um mesmo experimento, mas separados espacialmente por grandes dist�ncias, e mesmo
temporalmente, entraria em conflito com o conceito de observador na relatividade
especial, em particular com o conceito de simetrias do espa�otempo. Uma equa��o como na
representa��o de Schrodinger deve ser linear no tempo:
%
%
\[H\psi = i\hbar\deli{\psi}{t}\]
%
%
O artif�cio utilizado por Dirac para encontrar sua equa��o foi bastante engenhoso. Al�m
das hip�teses acima de mesclar conceitos da qu�ntica com a relatividade especial, Dirac
tamb�m imp�s que sua descri��o din�mica deveria ser uma equa��o diferencial linear no
tempo, na forma da equa��o de Schrodinger. A equa��o deveria ser Lorentz Covariante, e a
energia deveria ter a corre��o relativ�stica, $E^2 = (pc)^2 + (mc^2)^2$. Al�m disso,
deveria satisfazer tamb�m a Equa��o de Klein Gordon, o que imediatamente significa que
deve tamb�m satisfazer a equa��o de continuidade.
\\
Al�m disso, para que as simetrias requeridas pelo espa�otempo na relatividade especial
fossem satisfeitas, seria necess�rio que a equa��o fosse linear tamb�m no momento. Dirac
sabia que partir da energia total, $E = \sqrt{(pc)^2 + (mc^2)^2}$ e expandir em s�rie de
Taylor, resolvendo formalmente para termos de primeira ordem em p, n�o seria uma
boPartindoa sa�da, por isso resolveu fatorar a Equa��o de Klein Gordon. Uma fatora��o
simples, supondo solu��es escalares, seria (Observa-se que ser� usada a rela��o $a =
{\lambda_c}^{-1} = mc/\hbar$):
%
%
\[\left(\partial^2 - a^2\right)\psi = \left(\partial - a\right)\left(\partial +
a\right)\psi = 0\]
%
%
Como feito na equa��o de ondas (escalar). Entretanto, Dirac teve a id�ia de fazer essa
fatora��o utilizando operadores, deixando sua equa��o em uma forma geral:
%
%
\[
\left(\beta mc^2 + \versor{\alpha}\cdot \versor{p} \right)\psi =
i \hbar \deli{\psi}{t}
\]
%
%
Ele observou que os coeficientes $\alpha$ e $\beta$ n�o poderiam ser n�meros, mas sim
matrizes quadradas de ordem $n$ ainda desconhecida. Aplicando novamente o operador
(elevando ao quadrado sua equa��o):
%
%
\[
\left(\beta mc^2 + \versor{\alpha}\cdot \versor{p}c \right)
\left(\beta mc^2 + \versor{\alpha}\cdot \versor{p}c \right)\psi
= - \hbar^2 \deli{^2}{t^2}\psi
\]
%
%
Lembrando que estamos lidando com operadores, ao aplicar a distribui��o na multiplica��o,
tem-se o seguinte resultado:
%
%
\[
 (mc^2)^2\beta^2 + mc^3\left[\beta(\versor{\alpha}\cdot\versor{p}) +
(\versor{\alpha}\cdot\versor{p})\beta\right] + c^2(\versor{\alpha}\cdot\versor{p})^2
= - \hbar^2 \deli{^2}{t^2}\psi
\]
%
%
Na nota��o de Einstein de tensores, temos que:
%
%
\[
(mc^2)^2\beta^2 + mc^3(\beta \alpha_i p_i + \alpha_i p_i \beta) +
c^2(\alpha_i p_i)(\alpha_j p_j) = - \hbar^2 \deli{^2}{t^2}\psi
\]
%
%
As matrizes $\alpha_i$ e $\beta$ foram chamadas por Dirac de vari�veis
din�micas\index{vari�veis din�micas}, sendo elas independentes das vari�veis de momento,
$p_\mu$ o que implica que as vari�veis din�micas comutariam com as vari�veis $x_\mu$.
Lembrando que, $x_0 = t$. Utilizando a isotropia do espa�o, Dirac argumentou de forma
an�loga para as vari�veis $x_\mu$, o que implicaria que as vari�veis din�micas tamb�m
comutariam com as vari�veis de momento:
%
%
\[p_i\alpha_i = \alpha_i p_i \quad , \quad p_i\beta = \beta p_i\]
%
%
Podemos ainda fazer a seguinte manipula��o:
%
%
\[\alpha_i\alpha_j = \frac{1}{2}(\alpha_i\alpha_j + \alpha_j\alpha_i)\]
%
%
Sabendo desses resultados, encontra-se que:
%
%
\[
(mc^2)^2\beta^2 + mc^3(\beta\alpha_i + \alpha_i\beta)p_i +
\frac{c^2}{2}(\alpha_i\alpha_j + \alpha_j\alpha_i)p_i p_j = - \hbar^2 \deli{^2}{t^2}\psi
\]
%
%
Como a equa��o de Dirac deve satisfazer a equa��o de Klein Gordon, lembrando que a
equa��o de Klein Gordon � dada por (Na nota��o de Einstein, $\nabla^2 =
\delta_{ij}p_ip_j$):
%
%
\[ (- \delta_{ij}p_ip_j + m^2c^4)= - \hbar^2 \deli{^2}{t^2}\psi\]
%
%
De modo que, igualando os termos da equa��o de Klein Gordon, e da equa��o de Dirac
"elevada ao quadrado", tem-se que:
%
%
\[
(mc^2)^2\beta^2 + mc^3(\beta\alpha_i + \alpha_i\beta)p_i +
\frac{c^2}{2}(\alpha_i\alpha_j + \alpha_j\alpha_i)p_i p_j = - \delta_{ij}p_ip_j + (mc^2)^2
\]
%
%
Igualando os termos dos dois lados da equa��o, encontra-se as seguintes rela��es para as
vari�veis din�micas:
%
%
\[
\beta^2 = (\alpha_i)^2 = I
\quad , \quad
\alpha_i\alpha_j + \alpha_j\alpha_i = 2\delta_{ij}
\quad , \quad
\beta\alpha_i + \alpha_i\beta = 0
\]
%
%
Na pr�xima sess�o ser� explicado algumas sutilezas sobre as matrizes de Dirac, que em sua
representa��o, s�o representadas por:
%
%
\[    \beta = \begin{pmatrix} 1 & 0 & 0 & 0 \\ 0 & 1 & 0 & 0 \\ 0 & 0 & -1 & 0 \\
0 & 0
& 0 & -1 \end{pmatrix},\quad \alpha^1 \!=\! \begin{pmatrix} 0 & 0 & 0 & 1 \\ 0 & 0 & 1 & 0
\\ 0 & -1 & 0 & 0 \\ -1 & 0 & 0 & 0 \end{pmatrix}\]
%
%
\[\alpha^2 \!=\! \begin{pmatrix} 0 & 0 & 0 & -i \\ 0 & 0 & i & 0 \\ 0 & i & 0 & 0 \\ -i
& 0 & 0 & 0 \end{pmatrix},\quad \alpha^3 \!=\! \begin{pmatrix} 0 & 0 & 1 & 0 \\ 0 & 0 & 0
& -1 \\ -1 & 0 & 0 & 0 \\ 0 & 1 & 0 & 0 \end{pmatrix}\]
%
%
Ou de uma forma mais compacta:
%
%
\[
\alpha_k = \begin{pmatrix} 0 & \sigma_k \\ \sigma_k & 0 \end{pmatrix}
\qquad , \qquad
\beta = \begin{pmatrix} I_2 & 0 \\ 0 & -I_2 \end{pmatrix}
\]
%
%
Onde $\sigma_k$ $(k = 1,2,3)$ s�o as matrizes de Pauli, e $I_2$ � a identidade de
ordem 2.
\section{Grupo de Lorentz}

\subsection{Matrizes de Dirac}

As rela��es de anticomuta��o das vari�veis din�micas de Dira,
%
%
%
%
\[
\beta^2 = (\alpha_i)^2 = I
\quad , \quad
\alpha_i\alpha_j + \alpha_j\alpha_i = 2\delta_{ij}
\quad , \quad
\beta\alpha_i + \alpha_i\beta = 0
\]

Definem uma �lgebra para as matrizes $\beta$, $\alpha_i$, chamada de �lgebra de Dirac
\index{�lgebra de Dirac} que nada mais � do que um tipo de �lgebra de Clifford
\index{�lgebra de Clifford}, que � representada pelas matrizes $\gamma^\mu$:
%
%
\[    \gamma^0 = \begin{pmatrix} 1 & 0 & 0 & 0 \\ 0 & 1 & 0 & 0 \\ 0 & 0 & -1 & 0 \\
0 & 0
& 0 & -1 \end{pmatrix},\quad \gamma^1 \!=\! \begin{pmatrix} 0 & 0 & 0 & 1 \\ 0 & 0 & 1 & 0
\\ 0 & -1 & 0 & 0 \\ -1 & 0 & 0 & 0 \end{pmatrix}\]
%
%
\[\gamma^2 \!=\! \begin{pmatrix} 0 & 0 & 0 & -i \\ 0 & 0 & i & 0 \\ 0 & i & 0 & 0 \\ -i
& 0 & 0 & 0 \end{pmatrix},\quad \gamma^3 \!=\! \begin{pmatrix} 0 & 0 & 1 & 0 \\ 0 & 0 & 0
& -1 \\ -1 & 0 & 0 & 0 \\ 0 & 1 & 0 & 0 \end{pmatrix}\]

Esse conjunto de matrizes forma um espa�o pseudo ortogonal, de dimens�o 4. Um outro grupo
de matrizes que obedecem as mesmas regras de anticomuta��o s�o as matrizes de Pauli, que
s�o matrizes ortogonais em um espa�o de tr�s dimens�es, sendo portanto um subgrupo do
grupo das matrizes $\gamma^\mu$.

Al�m das propriedades de anticomuta��o, as vari�veis din�micas de Dirac apresentam outras
propriedades importantes:

\begin{enumerate}
\item As matrizes $\alpha_i$ e $\beta$ s�o hermitianas \index{hermitianas}. Este fato �
imediato da hermiticidade do Hamiltoniano. Fato este que implica que seus auto valores
sejam reais.
\item A rela��o $\beta^2 = (\alpha_i)^2 = I$ implica que estas matrizes possuem o
determinante igual � $\pm 1$. As matrizes cujo determinante � igual � $1$, pertencem ao
grupo das matrizes (de ordem n) ortogonais especiais\index{matrizes ortogonais especiais},
SO(n). Estas matrizes quando aplicadas como transforma��es lineares, atuam como rota��es.
\item Da propriedade $\beta\alpha_i + \alpha_i\beta = 0$, tem-se que $\alpha_i = -
\beta\alpha_i\beta$.
\item Pela invari�ncia do tra�o, e da propriedade de que $\traco(AB) = \traco(BA)$, e
utilizando a propriedade anterior, tem-se que:
%
%
\[\traco(\alpha_i) = \traco(-\beta\alpha_i\beta) = - \traco(\beta^2\alpha_i) = -
\traco(\alpha_i) \to \traco(\alpha_i) = 0\]
%
%
\item A rela��o $\beta^2 = (\alpha_i)^2 = I$ implica tamb�m que estas matrizes podem
possuir apenas auto valores iguais � $\pm 1$. A demonstra��o deste fato � imediata, pois
quando aplicada � um auto vetor, $\alpha_i v = av$, aplicando novamente a matriz, temos
que, $(\alpha_i)^2v = Iv = a^2v \to a = \pm 1$. A demonstra��o para $\beta$ � an�loga.
\item O fato de que as matrizes $\beta$ e $\alpha_i$ possuem auto valores apenas iguais �
$\pm 1$, implica que a ordem da matriz deve ser par, pois devem possuir a mesma propor��o
de auto valores negativos, quanto positivos. 
\end{enumerate}

As propriedades 5 e 6 nos remetem a uma discuss�o interessante, sobre o por qu� das
matrizes de Dirac serem de ordem 4. Como j� foi demonstrado que as matrizes devem possuir
uma ordem n = par, testemos a menor ordem poss�vel primeiramente, $n = 2$. � imposs�vel
determinar um par de matrizes que obede�am as propriedades anticomutativas da �lgebra de
Dirac. O �nico conjunto de matrizes que obedece essa �lgebra pertencem a um espa�o de 3
dimens�es, s�o as matrizes de Pauli. Esses fatos evidenciam que a menor ordem par
poss�vel para as matrizes seria $n = 4$.

Outras matrizes de Dirac podem ser obtidas atrav�s de transforma��es unit�rias. Isso
garante que as medidas f�sicas realizadas s�o independentes da representa��o utilizada,
na �lgebra de Dirac. Esse fato permite utilizar representa��es arbitrariamente de acordo
com a simplifica��o do problema f�sico em quest�o. Dirac usou como representa��o:
%
%
\[
\alpha_k = \begin{pmatrix} 0 & \sigma_k \\ \sigma_k & 0 \end{pmatrix}
\qquad , \qquad
\beta = \begin{pmatrix} I_2 & 0 \\ 0 & -I_2 \end{pmatrix}
\]
\section{Equa��o de Dirac}

\subsection{Part�cula Livre (Solu��o Geral)}

A equa��o de Dirac\index{equa��o de Dirac} � dada na forma covariante\index{covariante} da
Equa��o de onda na forma de Schrodinger\index{Schrodinger}:
%
%
\[H\psi = i\hbar\deli{\psi}{t}\]

Dirac escolheu que a equa��o fosse linear al�m do tempo, no momento tamb�m, de forma que:
%
%
\[\left(\versor{\alpha}\cdot pc + \versor{\beta} mc^2\right)\psi = i\hbar\deli{\psi}{t}\]

Onde $\alpha_i$, para $i = 1,2,3$, e $\beta$ s�o as matrizes de Dirac, o momento continua
tendo a mesma interpreta��o que na equa��o de Schrodinger $p \to i\hbar\nabla$. As
matrizes de Dirac apresentam a seguinte forma:
%
%
\[
\alpha_i = \begin{pmatrix}0 & \sigma_i\\ \sigma_i & 0\end{pmatrix} \qquad
\beta = \begin{pmatrix}I & 0\\ 0 & - I\end{pmatrix}
\]

Onde $I$ � a matriz identidade, e $\sigma_i$ s�o as matrizes de Pauli\index{matrizes de
Pauli}:
%
%
\[
\sigma_1 = \begin{pmatrix} 0 & 1 \\ 1 & 0 \end{pmatrix} \qquad
\sigma_2 = \begin{pmatrix} 0 & -i \\ i & 0 \end{pmatrix}  \qquad
\sigma_3 = \begin{pmatrix} 1 & 0 \\ 0 & -1 \end{pmatrix}
\]

A solu��o para a part�cula livre\index{part�cula livre} de Dirac possui como tentativa de
solu��o o seguinte Ansatz\index{Ansatz}:
%
%
\[\psi(\versor{x},t) = \psi(\versor{x})\exp(- i\epsilon t/\hbar)\]

Inserindo este Ansatz na equa��o de Dirac, encontra-se que:
%
%
\[H\psi(\versor{x}) = \epsilon\psi(\versor{x})\]

A auto solu��o $\epsilon$ possui como interpreta��o f�sica, ser a energia da part�cula
livre de Dirac, e possui valores tanto negativos como positivos, o que resulta em uma
nova interpreta��o, muito diferente daquela da resolu��o da equa��o de Schrodinger.
Deve-se observar tamb�m que $\epsilon$ faz parte da solu��o temporal do
spinor\index{spinor} $\psi(\versor{x},t)$, sendo respons�vel pela evolu��o no tempo, para
a solu��o estacion�ria, $\psi(\versor{x})$.

Nos adiantando a esta an�lise dos valores negativos e positivos de energia\index{energia},
ser� feita uma proposta de separar o spinor da seguinte maneira:
%
%
\[
\psi_+ = \begin{pmatrix}\psi_1 \\ \psi_2\end{pmatrix} \qquad
\psi_- = \begin{pmatrix}\psi_3 \\ \psi_4\end{pmatrix} \quad \to \quad
\psi(\versor{x}) = \begin{pmatrix}\psi_+ \\ \psi_-\end{pmatrix}\]

Representando a equa��o de Dirac em forma matricial, tem-se que:
%
%
\[
\begin{pmatrix}0 & \versor\sigma\cdot \versor{p} c \\ \versor\sigma\cdot \versor{p} c &
0\end{pmatrix}
\begin{pmatrix}\psi_+ \\ \psi_-\end{pmatrix} +
mc^2\begin{pmatrix}I & 0\\ 0 & - I\end{pmatrix}
\begin{pmatrix}\psi_+ \\ \psi_-\end{pmatrix} =
\epsilon\begin{pmatrix}\psi_+ \\ \psi_-\end{pmatrix}
\]

Tem-se ainda que, ap�s inserir as formas das matrizes de Pauli, e da Identidade:
%
%
\[
\begin{pmatrix}
mc^2          & 0             & p_3c          & (p_1 - ip_2)c \\
0             & mc^2          & (p_1 + ip_2)c & - p_3c \\
p_3c          & (p_1 - ip_2)c & mc^2          & 0 \\
(p_1 + ip_2)c & - p_3c        & 0             & mc^2
\end{pmatrix}
\begin{pmatrix}\psi_1 \\ \psi_2 \\ \psi_3 \\ \psi_4 \end{pmatrix} =
\epsilon\begin{pmatrix}\psi_1 \\ \psi_2 \\ \psi_3 \\ \psi_4 \end{pmatrix}
\]

O sistema linear acima, corresponde �s seguintes equa��es, nas vari�veis, $\psi_i$, sendo
$i = 1,2,3,4$:
%
%
\[\begin{array}{ccccccccc}
(\epsilon - mc^2)\psi_1 &+& 0\psi_2 &-& p_3c\psi_3 &-& c(p_1 - ip_2)\psi_4    &=& 0\\
0\psi_1 &+& (\epsilon - mc^2)\psi_2 &-& c(p_1 + ip_2)\psi_3 &+& p_3c\psi_4  &=& 0\\
- p_3c\psi_1 &-& c(p_1 - ip_2)\psi_2 &+& (\epsilon - mc^2)\psi3 &+& 0\psi_4 &=& 0\\
-c(p_1 - ip_2)\psi_1 &+& p_3c\psi_2 &+& 0\psi_3 &+& (\epsilon - mc^2)\psi_4 &=& 0
\end{array}
\]

O conjunto de equa��es acima � homog�neo\index{homog�neo}, logo � poss�vel e determinado,
desde que seu determinante seja nulo, ou seja, o equivalente a resolver o sistema linear,
$H\psi = \epsilon\psi$, de modo que, $|H - I\epsilon| |\psi| = 0$. O determinante
\index{determinante} (polin�mio caracter�stico) possui como resultado (ra�zes), os
seguintes valores para o auto valor, $\epsilon$:
%
%
\[\epsilon_+ = \sqrt{mc^2 + (pc)^2} \quad , \quad \epsilon_- = - \sqrt{mc^2 + (pc)^2}\]

Para o auto valor\index{auto valor}, $\epsilon_+$, existem duas solu��es linearmente
independentes\index{linearmente independentes}.
%
%
\[\psi_1^{(1)} = 1
\ \ , \ \
\psi_2^{(1)} = 0
\ \ , \ \
\psi_3^{(1)} = \frac{p_3c}{\epsilon_+ + mc^2}
\ \ , \ \
\psi_4^{(1)} = \frac{(p_1 + ip_2)c}{\epsilon_+ + mc^2}
\]
%
%
\[
\psi_1^{(2)} = 0
\ \ , \ \
\psi_2^{(2)} = 1
\ \ , \ \
\psi_3^{(2)} = \frac{(p_1 - ip_2)c}{\epsilon_+ + mc^2}
\ \ , \ \
\psi_4^{(2)} = - \frac{p_3c}{\epsilon_+ + mc^2}
\]

Da mesma forma, temos que, para o auto valor, $\epsilon_-$, existem duas solu��es
linearmente independentes.
%
%
\[\psi_1^{(1)} = \frac{p_3c}{\epsilon_- - mc^2}
\ \ , \ \ 
\psi_2^{(1)} = \frac{(p_1 + ip_2)c}{\epsilon_- - mc^2}
\ \ , \ \
\psi_3^{(1)} = 1
\ \ , \ \
\psi_4^{(1)} = 0 
\]
%
%
\[\psi_1^{(2)} = \frac{(p_1 - ip_2)}{\epsilon_- - mc^2}
\ \ , \ \
\psi_2^{(2)} = - \frac{p_3c}{\epsilon_- - mc^2}
\ \ , \ \
\psi_3^{(2)} = 0
\ \ , \ \
\psi_4^{(2)} = 1
\]


\subsection{Part�cula Livre (Solu��o Mais Simples)}

A solu��o mais simples para a part�cula livre\index{part�cula livre} de Dirac, � obtida
quando o vetor momento, $\versor{p}$, apontar na dire��o positiva do eixo $z$, nesse caso,
a equa��o de Dirac, j� em sua forma matricial, toma a seguinte apar�ncia:
%
%
\[
\begin{pmatrix}
mc^2 & 0      & pc  & 0 \\
0    & mc^2   & 0     & - pc \\
pc & 0      & mc^2  & 0 \\
0    & - pc & 0     & mc^2
\end{pmatrix}
\begin{pmatrix}\psi_1 \\ \psi_2 \\ \psi_3 \\ \psi_4 \end{pmatrix} =
\epsilon\begin{pmatrix}\psi_1 \\ \psi_2 \\ \psi_3 \\ \psi_4 \end{pmatrix}
\]

Repare que por simplicidade de nota��o, foi definido que, $p_3 = p$ .De forma que se tenha
o seguinte conjunto de equa��es homog�neas nas vari�veis, $\psi_i$, sendo, $i = 1,2,3,4$:
%
%
\[\begin{array}{ccccccccc}
(\epsilon - mc^2)\psi_1 &+& 0\psi_2 &-& pc\psi_3 &-& 0\psi_4    &=& 0\\
0\psi_1 &+& (\epsilon - mc^2)\psi_2 &-& 0\psi_3 &+& pc\psi_4  &=& 0\\
- pc\psi_1 &-& \psi_2 &+& (\epsilon - mc^2)\psi3 &+& 0\psi_4 &=& 0\\
0\psi_1 &+& pc\psi_2 &+& 0\psi_3 &+& (\epsilon - mc^2)\psi_4 &=& 0
\end{array}
\]

O determinante\index{determinante} (polin�mio caracter�stico) possui como resultado
(ra�zes), os seguintes
valores para o auto valor, $\epsilon$:
%
%
\[\epsilon_+ = \sqrt{mc^2 + (pc)^2} \quad , \quad \epsilon_- = - \sqrt{mc^2 + (pc)^2}\]

Para o auto valor\index{auto valor}, $\epsilon_+$, existem duas solu��es linearmente
independentes.
%
%
\[\psi_1^{(1)} = 1
\ \ , \ \
\psi_2^{(1)} = 0
\ \ , \ \
\psi_3^{(1)} = \frac{pc}{\epsilon_+ + mc^2}
\ \ , \ \
\psi_4^{(1)} = 0
\]
%
%
\[
\psi_1^{(2)} = 0
\ \ , \ \
\psi_2^{(2)} = 1
\ \ , \ \
\psi_3^{(2)} = 0
\ \ , \ \
\psi_4^{(2)} = - \frac{pc}{\epsilon_+ + mc^2}
\]

Da mesma forma, temos que, para o auto valor, $\epsilon_-$, existem duas solu��es
linearmente independentes.
%
%
\[\psi_1^{(1)} = \frac{pc}{\epsilon_- - mc^2}
\ \ , \ \
\psi_2^{(1)} = 0
\ \ , \ \
\psi_3^{(1)} = 1
\ \ , \ \
\psi_4^{(1)} = 0
\]
%
%
\[\psi_1^{(2)} = 0
\ \ , \ \
\psi_2^{(2)} = - \frac{pc}{\epsilon_- - mc^2}
\ \ , \ \
\psi_3^{(2)} = 0
\ \ , \ \
\psi_4^{(2)} = 1
\]

\subsection{Problema de Difus�o}

Seja uma barreira de potencial \index{barreira de potencial } quadrada, de altura $V_0 >
0$ e largura L, definida no intervalo $[0,L]$. Considerando o movimento de pacotes de onda
na dire��o z, de forma a ter-se tr�s dom�nios: 1 para $z\leq 0 $, 2 para $0 \leq z \leq
L$, e 3 para $z \geq L$. 
\\
Dessa maneira, tem-se que a equa��o de Dirac para as regi�es 1 e 3 � dada por,
$(\versor{\alpha} \cdot \versor{p} c + \beta mc^2)\psi = E\psi$, enquanto que para a
regi�o 2, a equa��o de Dirac � dada pela seguinte rela��o:
%
%
\[(\versor{\alpha} \cdot \versor{p} c + \beta mc^2)\psi = (E - V_0)\psi\]
%
%
As solu��es (spinores) \index{spinores} da equa��o de Dirac nas regi�es 1 e 3 s�o as
solu��es da part�cula livre \index{part�cula livre} de Dirac, e como a escolha
unidirecional foi na dire��o z, os spinores s�o fun��es apenas desta coordenada, ou seja,
$\psi = \psi(z)$.
\\
Para a regi�o 2, as solu��es s�o as mesmas, por�m com a seguinte
substitui��o, $E \to E - V_0$. Sendo que o espectro de energias pode assumir qualquer
valor no intervalo aberto $(-\infty,\infty)$, descrevendo dessa maneira al�m de
part�culas, as chamadas anti-part�culas.
\\
As solu��es da equa��o de Dirac admitem spins "up" e "down", mas o tratamento ser� dado
sem o spin flip \index{spin flip}, o que torna a escolha arbitr�ria da orienta��o de spin,
neste caso o spin up \index{spin up}. Abaixo as solu��es para as respectivas regi�es:
%
%
\begin{enumerate}
 \item Onde $(z \leq 0)$, sendo $\psi_1(z)$ o spinor para essa regi�o, e sendo $p$ o
 valor do momento, e seja $\epsilon_1 = E + mc^2$, sendo respeitada a seguinte
 rela��o, $E^2 = (pc)^2 + (mc^2)^2$ tem-se que que:
 %
 %
 \[\psi_1(z) = A_1\exp(ipz/\hbar)\begin{pmatrix} 1 \\ 0 \\ pc/\epsilon_1 \\
 0\end{pmatrix} + A_2\exp(- ipz/\hbar)\begin{pmatrix}1 \\ 0 \\ - pc/\epsilon_1 \\
 0\end{pmatrix}\]
 %
 %
 \item Onde $(0 \leq z \leq L)$, sendo $\psi_2(z)$ o spinor para essa regi�o, e sendo $q$
 o valor do momento, e seja $\epsilon_2 = E - V_0 + mc^2$, sendo respeitada a seguinte
 rela��o, $(E - V_0)^2 = (qc)^2 + (mc^2)^2$ tem-se que que:
 %
 %
 \[\psi_2(z) = B_1\exp(iqz/\hbar)\begin{pmatrix} 1 \\ 0 \\ qc/\epsilon_2 \\
 0\end{pmatrix} + B_2\exp(- iqz/\hbar)\begin{pmatrix}1 \\ 0 \\ - qc/\epsilon_2 \\
 0\end{pmatrix}\]
 %
 %
 \item Onde $(z \geq L)$, sendo $\psi_3(z)$ o spinor para essa regi�o, e sendo
 $p$ o valor do momento, e seja $\epsilon_1 = E + mc^2$, sendo respeitada a seguinte
 rela��o, $E^2 = (pc)^2 + (mc^2)^2$ tem-se que que:
 %
 %
 \[\psi_3(z) = C_1\exp(ipz/\hbar)\begin{pmatrix} 1 \\ 0 \\ pc/\epsilon_1 \\
 0\end{pmatrix} + C_2\exp(- ipz/\hbar)\begin{pmatrix}1 \\ 0 \\ - pc/\epsilon_1 \\
 0\end{pmatrix}\]
\end{enumerate}
%
%
Para que a densidade de corrente \index{densidade de corrente} $(\partial_\mu j^\mu)$ seja
conservada � necess�rio que hajam condi��es de continuidade das fun��es de onda, nas
fronteiras da barreira de potencial. Deve-se ter ent�o que $\psi_1(0) = \psi_2(0)$ e que
$\psi_2(L) = \psi_3(L)$. Fazendo essas compara��es, encontra-se que:
%
%
\begin{eqnarray}
\nonumber \psi_1(0) &=& \psi_2(0)\\
\nonumber  A_1\begin{pmatrix} 1 \\ 0 \\ pc/\epsilon_1 \\  0\end{pmatrix} +
A_2\begin{pmatrix}1 \\ 0 \\ - pc/\epsilon_1 \\ 0\end{pmatrix} &=& B_1\begin{pmatrix} 1 \\
0 \\ qc/\epsilon_2 \\ 0\end{pmatrix} + B_2\begin{pmatrix}1 \\ 0 \\ - qc/\epsilon_2 \\
0\end{pmatrix}
\end{eqnarray}
%
%
Dando resultado ao seguinte par de equa��es, nas vari�veis $A_1, A_2, B_1, B_2$:
%
%
\begin{eqnarray}
\nonumber A_1 + A_2 &=& B_1 + B_2\\
\nonumber (A_1 - A_2)\frac{pc}{\epsilon_1} &=& (B_1 - B_2)\frac{qc}{\epsilon_2}
\end{eqnarray}
%
%
� poss�vel reescrever o sistema de equa��es acima em forma matricial:
%
%
\[\begin{pmatrix}1 & 1 \\ 1 & -1\end{pmatrix}\begin{pmatrix}A_1 \\ A_2\end{pmatrix} =
\begin{pmatrix}1 & 1 \\ q\epsilon_1/p\epsilon_2 & -
q\epsilon_1/p\epsilon_2\end{pmatrix}\begin{pmatrix}B_1 \\ B_2\end{pmatrix}\]
%
%
De modo que (utilizando a seguinte simplifica��o, $\alpha = q\epsilon_1/p\epsilon_2$):
%
%
\begin{eqnarray}
\nonumber \begin{pmatrix}A_1 \\ A_2\end{pmatrix} &=& \begin{pmatrix}1/2 & 1/2 \\ 1/2 &
-1/2\end{pmatrix}\begin{pmatrix}1 & 1 \\ \alpha & -\alpha\end{pmatrix}\begin{pmatrix}B_1
\\ B_2\end{pmatrix}\\
\nonumber &=& \frac{1}{2}\begin{pmatrix}1 + \alpha & 1 - \alpha \\ 1 - \alpha &
1 + \alpha\end{pmatrix} \begin{pmatrix}B_1 \\ B_2\end{pmatrix}
\end{eqnarray}
%
%
Procedendo da mesma maneira para $z = L$, tem-se que:
%
%
\begin{eqnarray}
\nonumber \psi_2(L) &=& \psi_3(L)\\
\nonumber B_1e^{iqL/\hbar}\begin{pmatrix} 1 \\ 0 \\ qc/\epsilon_2 \\
 0\end{pmatrix} + B_2 e^{- iqL/\hbar}\begin{pmatrix}1 \\ 0 \\ - qc/\epsilon_2 \\
 0\end{pmatrix} &=& C_1e^{ipL/\hbar}\begin{pmatrix} 1 \\ 0 \\ pc/\epsilon_1 \\
 0\end{pmatrix} \\
\nonumber &+& C_2 e^{- ipL/\hbar}\begin{pmatrix}1 \\ 0 \\ - pc/\epsilon_1 \\
 0\end{pmatrix}
\end{eqnarray}
%
%
Dando resultado ao seguinte par de equa��es, nas vari�veis $B_1, B_2, C_1, C_2$:
%
%
\begin{eqnarray}
\nonumber  B_1e^{iqL/\hbar} + B_2e^{- iqL/\hbar} &=& C_1e^{ipL/\hbar} +
C_2e^{-ipL/\hbar}\\
\nonumber  \left(B_1e^{iqL/\hbar} - B_2e^{- iqL/\hbar}\right)\frac{qc}{\epsilon_2} &=&
\frac{pc}{\epsilon_1}\left(C_1e^{ipL/\hbar} - C_2e^{-ipL/\hbar}\right)
\end{eqnarray}
%
%
� poss�vel reescrever o sistema de equa��es acima em forma matricial. Para simplificar a
nota��o, ser� feito o seguinte apoio alg�brico, $\beta = e^{iqL/\hbar}$, de modo que,
$\beta^* = e^{- iqL/\hbar}$, e $\beta\beta^* = 1$, da mesma maneira ser� definido que,
$\gamma = e^{ipL/\hbar}$, e que $\gamma^* = e^{- iqL/\hbar}$, sendo, $\gamma\gamma^* =
1$. Al�m disso, deve-se lembrar que, $\alpha = q\epsilon_1/p\epsilon_2$
%
%
\[\begin{pmatrix}\beta & \beta^* \\ \beta & - \beta^*\end{pmatrix}\begin{pmatrix}B_1
\\ B_2\end{pmatrix} =
\begin{pmatrix}\gamma & \gamma^* \\ \gamma/\alpha &
- \gamma^*/\alpha\end{pmatrix}\begin{pmatrix}C_1 \\ C_2\end{pmatrix}\]
%
%
De modo que:
%
%
\begin{eqnarray}
\nonumber \begin{pmatrix}B_1 \\ B_2\end{pmatrix} &=& \frac{1}{2}\begin{pmatrix}1/\beta &
1/\beta \\ 1/\beta^* & - 1/\beta^*\end{pmatrix}\begin{pmatrix}\gamma & \gamma^* \\
\gamma/\alpha & - \gamma^*/\alpha\end{pmatrix}\begin{pmatrix}C_1 \\ C_2\end{pmatrix}\\
\nonumber &=& \frac{1}{2}\begin{pmatrix}\frac{\gamma}{\beta}\frac{1 + \alpha}{\alpha} &
\frac{\gamma^*}{\beta}\frac{\alpha - 1}{\alpha}\\
\frac{\gamma}{\beta^*}\frac{\alpha - 1}{\alpha} &
\frac{\gamma^*}{\beta^*}\frac{1 + \alpha}{\alpha}\\
\end{pmatrix}\begin{pmatrix}C_1 \\ C_2\end{pmatrix}
\end{eqnarray}
%
%
De forma que finalmente se tenha:
%
%
\begin{eqnarray}
\nonumber \begin{pmatrix}A_1 \\ A_2\end{pmatrix} &=& \frac{1}{4\alpha}
\begin{pmatrix}1 + \alpha & 1 - \alpha\\
 1 - \alpha & 1 + \alpha
\end{pmatrix}
\begin{pmatrix}\frac{\gamma}{\beta}(1 + \alpha) & \frac{\gamma^*}{\beta}(\alpha - 1)\\
\frac{\gamma}{\beta^*}(\alpha - 1)& \frac{\gamma^*}{\beta^*}(1 + \alpha)\\
\end{pmatrix}
\begin{pmatrix}
C_1 \\
C_2
\end{pmatrix}\\
\nonumber &=& \begin{pmatrix}
\frac{\gamma}{\beta}(1 + \alpha)^2 - \frac{\gamma}{\beta^*}(\alpha - 1)^2 &
(\alpha^2 - 1)\left(\frac{\gamma^*}{\beta} - \frac{\gamma^*}{\beta^*}\right)\\
(1 - \alpha^2)\left(\frac{\gamma^*}{\beta} - \frac{\gamma^*}{\beta^*}\right) &
\frac{\gamma}{\beta^*}(1 + \alpha)^2 - \frac{\gamma}{\beta}(\alpha - 1)^2
\end{pmatrix}\begin{pmatrix}
C_1 \\
C_2
\end{pmatrix}
\end{eqnarray}

\newpage\noindent
Lembrando que $\beta = e^{iqL/\hbar}$, $\beta^* = e^{- iqL/\hbar}$, $\beta\beta^* = 1$,
$\gamma = e^{ipL/\hbar}$, $\gamma^* = e^{- iqL/\hbar}$, e que $\alpha =
q\epsilon_1/p\epsilon_2$, tem-se que:
%
%
\begin{multicols}{2}
\begin{enumerate}
\item $\ds \frac{\gamma}{\beta}     = \exp\left(iL(p - q)/\hbar\right)$
\item $\ds \frac{\gamma}{\beta^*}   = \exp\left(iL(p + q)/\hbar\right)$
\item $\ds \frac{\gamma^*}{\beta}   = \exp\left(- iL(p + q)/\hbar\right)$
\item $\ds \frac{\gamma^*}{\beta^*} = \exp\left(- iL(p - q)/\hbar\right)$
\end{enumerate}
\end{multicols}
%
%
Lembrando que $\epsilon_1 = E + mc^2$, e que $\epsilon_2 = E - V_0 + mc^2$, sendo que
temos as seguintes rela��es entre os momentos, em fun��o de $E, V_0, mc^2$:
%
%
\begin{eqnarray}
\nonumber p &=& \pm \frac{1}{c}\sqrt{E^2 - (mc^2)^2}\\
\nonumber q &=& \pm \frac{1}{c}\sqrt{(E - V_0)^2 - (mc^2)^2}
\end{eqnarray}
%
%
De forma que a constante, $\alpha = q\epsilon_1/p\epsilon_2$, seja dada pela seguinte
rela��o:
%
%
\begin{eqnarray}
\nonumber \alpha &=& \frac{\epsilon_1}{\epsilon_2}\frac{q}{p}\\
\nonumber        &=& \frac{E + mc^2}{E - V_0 + mc^2}\sqrt{\frac{(E - V_0)^2 -
(mc^2)^2}{E^2 - (mc^2)^2}}\\
\nonumber        &=& \sqrt{\frac{(E + mc^2)(E - V_0 - mc^2)}{(E - mc^2)(E - V_0 + mc^2)}}
\end{eqnarray}


\begin{thebibliography}{}

\bibitem{DeLeoI}
De Leo, S. ; Rotelli, P. Above barrier Dirac multiple scattering and resonances. European Physical Journal C, v. 46, p. 551-558, 2006.

\bibitem{DeLeoII}
De Leo, S. ; Rotelli, P. Dirac equation studies in the tunnellinf energy zone. European Physical Journal C, v. 51, p. 241-247, 2007.

\bibitem{Dirac}
Dirac, P.A.M. The Principles of Quantum Mechanics, Oxford University Press, 1982.

\bibitem{Gottfried}
Gottfried, K. Quantum Mechanics: Fundamentals, Springer, 2003.

\bibitem{Greiner}
Greiner, W. Relativistic Quantum Mechanics, Wave Equations, Springer, 2000.

\bibitem{De Leo}
De Leo, S. ; Rotelli, P. . Above barrier Dirac multiple scattering and resonances. European Physical Journal C, v. 46, p. 551-558, 2006.

\end{thebibliography}

\printindex
\end{document}